% !TEX TS-program = xelatex
% !TEX encoding = UTF-8 Unicode
% !Mode:: "TeX:UTF-8"

\documentclass{resume}
\usepackage{zh_CN-Adobefonts_external} % Simplified Chinese Support using external fonts (./fonts/zh_CN-Adobe/)
%\usepackage{zh_CN-Adobefonts_internal} % Simplified Chinese Support using system fonts
\usepackage{linespacing_fix} % disable extra space before next section
\usepackage{cite}

\begin{document}
\pagenumbering{gobble} % suppress displaying page number

\name{徐永浩}

% {E-mail}{mobilephone}{homepage}
% be careful of _ in emaill address
\contactInfo{(+86) 181-6240-9012}{yonghao.xu@outlook.com}{游戏服务端开发工程师}{}
% {E-mail}{mobilephone}
% keep the last empty braces!
%\contactInfo{xxx@yuanbin.me}{(+86) 131-221-87xxx}{}
 
% \section{个人总结}
% 本人在校成绩优秀、乐观向上，工作负责、自我驱动力强、热爱尝试新事物，认同开放、连接、共享的Web在未来的不可替代性。在校期间长期从事可视分析(Web的2D/3D时空可视化)相关研究，对Web技术发展趋势及前端工程化解决方案有浓厚兴趣。\textbf{现任职于 BAT 集团。}

% \section{\faGraduationCap\ 教育背景}
\section{教育背景}
\datedsubsection{\textbf{上海交通大学},仪器仪表工程,\textit{硕士}}{2016.9 - 2019.6}
% \ \textbf{排名11/133(前10\%)},中国科学院大学学业奖学金(2次),IEEE Student member,预计2018年6月毕业
\datedsubsection{\textbf{上海交通大学},测控技术与仪器,\textit{本科}}{2012.9 - 2016.6}
% \ \textbf{排名2/62(前5\%)},国家励志奖学金,人民奖学金(7次),科技竞赛奖(2次),北京市普通高等学校优秀毕业生,北京理工大学优秀毕业生,软件学院金牌毕业生,优秀团员/优秀学生(5次)


% \section{\faCogs\ IT 技能}
\section{技术能力}
% increase linespacing [parsep=0.5ex]
\begin{itemize}[parsep=0.2ex]
  \item 掌握C、C++、Lua、Python编程语言
  \item 熟悉MongoDB、Etcd
  \item 熟悉游戏服务器架构，熟悉skynet框架
\end{itemize}

% \end{itemize}

\section{项目经历}

\datedsubsection{\textbf{英雄互娱 | 二重螺旋}, 服务端负责人}{2021-至今}
\begin{itemize}
  \item 基于skynet，从0到1搭建服务器框架，设计并实现核心功能模块
  \item 带领团队完成三轮测试和公测，服务器运行稳定，首日国服最高40w在线，300w注册
  \item 分布式架构，通过Etcd实现服务注册与发现，支持动态扩容
  \item 属性系统
  \item 压测与优化：实现压测系统，业务逻辑层优化（单点压力），数据库优化（分片/索引策略/读写分离/慢日志优化等），内存泄漏检测与内存优化（主要针对玩家数据），提高登录并发量，确定单服承载，提升服务器性能
  \item 审核代码，提出改进优化意见
  \item 线上服务器部署与维护
\end{itemize}

\datedsubsection{\textbf{网易 | 幻书启世录}, 服务端开发工程师}{2019-2021}
\begin{itemize}
  \item 11
\end{itemize}

% \begin{onehalfspacing}
% \end{onehalfspacing}

% \datedsubsection{\textbf{DID-ACTE} 荷兰莱顿}{2015年}
% \role{本科毕业设计}{LIACS 交换生}
% 利用结巴分词对中国古文进行分词与词性标注，用已有领域知识训练形成 classifier 并对结果进行调优
% \begin{onehalfspacing}
% \begin{itemize}
%   \item 利用结巴分词对中国古文进行分词与词性标注
%   \item 利用已有领域知识训练形成 classifier, 并用分词结果进行测试反馈
%   \item 尝试不同规则，对 classifier 进行调优
% \end{itemize}
% \end{onehalfspacing}

% \section{竞赛获奖/项目作品}
% % increase linespacing [parsep=0.5ex]
% \begin{itemize}[parsep=0.2ex]
% %   \item LeetCodeOJ Solutions, \textit{https://github.com/hijiangtao/LeetCodeOJ}
%   \item 第三届中国软件杯大学生软件设计大赛\textbf{全国一等奖}( \textit{http://www.cnsoftbei.com/} ),2014 年8月
%   \item 中国机器人大赛创意设计大赛\textbf{全国特等奖}( \textit{http://www.rcccaa.org/} ),2013年8月
% %   \item 中国机器人大赛暨Robocup公开赛(武术擂台赛)全国一等奖,2013年10月
%   \item 第11届北京理工大学“世纪杯”竞赛学生课外科技作品竞赛\textbf{特等奖},2013年8月
%   \item VIS Components for security system, \textit{https://hijiangtao.github.io/ss-vis-component/}
%   \item 个人博客：\textit{https://hijiangtao.github.io/}，更多作品见 \textit{https://github.com/hijiangtao}
% %   \item 电视节目"爸爸去哪儿"可视化分析展示, \textit{https://hijiangtao.github.io/variety-show-hot-spot-vis/}
% \end{itemize}

% \section{\faHeartO\ 项目/作品摘要}
% \section{项目/作品摘要}
% \datedline{\textit{An Integrated Version of Security Monitor Vis System}, https://hijiangtao.github.io/ss-vis-component/ }{}
% \datedline{\textit{Dark-Tech}, https://github.com/hijiangtao/dark-tech/ }{}
% \datedline{\textit{融合社交网络数据挖掘的电视节目可视化分析系统}, https://hijiangtao.github.io/variety-show-hot-spot-vis/}{}
% \datedline{\textit{LeetCodeOJ Solutions}, https://github.com/hijiangtao/LeetCodeOJ}{}
% \datedline{\textit{Info-Vis}, https://github.com/ISCAS-VIS/infovis-ucas}{}


% % \section{\faInfo\ 社会实践/其他}
% \section{社区参与/实践其他}
% % increase linespacing [parsep=0.5ex]
% \begin{itemize}[parsep=0.2ex]
%   \item 乐于参与开源社区讨论,\textbf{参与翻译 Vue.js, webpack, WebAssembly, Babel 文档，印记中文成员}
%   \item 中国科学院大学2016秋季学期可视化与可视分析课程助教，\textit{http://vis.ios.ac.cn/infovis-ucas/}
%   \item 未来论坛学生会成员、北理社联新闻信息中心主任、北理工软件学院学生会宣传部副部长(2012-2016)
%   \item 2013-2015 北京市共青团“温暖衣冬”志愿者,第九届园博会志愿者,2014 FLL机器人世锦赛志愿者
% \end{itemize}

%% Reference
%\newpage
%\bibliographystyle{IEEETran}
%\bibliography{mycite}
\end{document}
